\chapter{The Theodore Transform: Inter-Axis Projection and Cross-Domain Transposition}
\label{apx:theodore_transform}

This appendix names and formalizes an operation the manuscript has used throughout without naming
it. The operation transfers information from one domain to the next, the way a Laplace transform
carries a function from the time domain to the frequency domain: the object changes its
representation, and the structure the representation carries stays fixed. Chapter~\ref{ch:gendered}
uses the operation to project the Haitian Theorem from the racial axis onto the gendered axis.
Chapters~\ref{ch:redefining}--\ref{ch:algorithmic_epoch} use the operation to transpose the entire
social-mechanical framework into electrodynamics. This appendix defines the operator, states its
invariance condition, works both uses on material already in the manuscript, and marks the limit of
what the operator can establish.

\medskip

\noindent\textbf{Confidence tier.} The formal apparatus of this appendix is Tier~3 (structural).
It names and unifies moves the manuscript already executes and introduces no new calibrated
measurement. Falsification criteria appear in Section~\ref{sec:tt_falsification}.

\section{Domains, Vocabularies, and Correspondence Maps}
\label{sec:tt_domains}

Fix an indexed family of domains $\{D_\alpha\}_{\alpha \in A}$. Two index sets matter in this book.
The first is the set of identity axes along which the Predatory Min-Max Function partitions the
population: race, gender, sexuality, and the further axes Chapter~\ref{ch:spectral_carrier}
resolves into sub-bands. The second is the set of formal substrates across which the framework has
been developed: social mechanics, circuit theory, thermodynamics, and control systems
(Appendix~\ref{app:universality}).

Each domain $D_\alpha$ carries a vocabulary $\mathcal{V}_\alpha$ (the symbols and defined
quantities of that domain) and a claim set $\mathcal{C}_\alpha$ (the propositions statable in that
vocabulary). On the racial axis the vocabulary includes $O_{\text{racialized}}$,
$F_{\text{enforce}}$, the wage $\psi$, and the Enclosure Score $\mathcal{S}_{\text{enc}}$. On the
gendered axis the vocabulary includes $O_{\text{gendered}}$
(Equation~\ref{eq:12.7-gendered-outgroup-set}), coverture, and the reproductive extraction kernel
documented in Chapter~\ref{ch:gendered}.

\begin{definition}[Structural correspondence]
A \textbf{structural correspondence} between domains $D_\alpha$ and $D_\beta$ is a map
$\sigma_{\alpha\beta} : \mathcal{V}_\alpha \to \mathcal{V}_\beta$ assigning each element of the
source vocabulary its structural analogue in the target vocabulary: the element occupying the same
position in the extraction topology $(E, P_{\text{uppet}}, F_{\text{enforce}}, I_{\text{buffer}},
O)$ and standing in the same relations to its neighbors.
\end{definition}

The manuscript constructs such correspondences repeatedly and explicitly.
$O_{\text{racialized}} \mapsto O_{\text{gendered}}$ maps the racialized Out-group to the gendered
Out-group; both sets are defined by exclusion from lethal autonomy
(Equation~\ref{eq:6.1-lethal-autonomy-gradient}). The slave-patrol and carceral machinery of
Chapter~\ref{ch:enforcement} maps to the gendered enforcement machinery of
Chapter~\ref{ch:gendered}: coverture statutes, eugenics law, and the criminalization of pregnancy.
The kinetic labor of $O_{\text{racialized}}$ and $I_{\text{buffer}}$ maps to the power supply
$V_{cc}$ of the circuit model, and the psychological wage $\psi_s$ maps to reactive power
$Q_{\text{reactive}}$ (Appendix~\ref{apx:extraction_chart}).

\section{The Transform}
\label{sec:tt_transform}

\begin{definition}[The Theodore Transform]
Let $\sigma_{\alpha\beta}$ be a structural correspondence. The \textbf{Theodore Transform} is the
substitution operator
\begin{equation}\label{eq:tt.1-transform}
\TT_{\alpha \to \beta} \;:\; \mathcal{C}_\alpha \;\longrightarrow\; \mathcal{C}_\beta,
\qquad c \;\longmapsto\; \sigma_{\alpha\beta}(c),
\end{equation}
where $\sigma_{\alpha\beta}(c)$ denotes the claim obtained by replacing every vocabulary element
of $c$ with its image under $\sigma_{\alpha\beta}$.
\end{definition}

The operator carries a claim across domains the way the Laplace transform carries a function
across representations. The input and the image differ in vocabulary. The image is legible exactly
when the correspondence preserves the relations the claim asserts. The symbol $\TT$ is
reserved for this operator throughout the manuscript, consistent with the notation rule of
Chapter~\ref{ch:system_init}: no other quantity in this book carries it.

\section{The Invariance Condition}
\label{sec:tt_invariance}

A claim asserted as a principle claims scope. A claimant who asserts a normative claim $c$ on
domain $D_\alpha$ as a principle commits to the transformed claim
$\TT_{\alpha \to \beta}(c)$ on every domain $D_\beta$ inside the principle's scope. A
claimant who asserts anti-racism as a principle thereby asserts the image of that claim on the
gendered axis, and on every other axis the principle covers. The transform converts that commitment
into a checkable condition:

\begin{equation}\label{eq:tt.2-invariance}
c \text{ holds on } D_\alpha \quad\Longrightarrow\quad
\TT_{\alpha \to \beta}(c) \text{ holds on } D_\beta
\qquad \forall\, \beta \in S(c),
\end{equation}

where $S(c)$ is the scope the claimant asserts. When Equation~\ref{eq:tt.2-invariance} holds across
$S(c)$, the claim is applied \textbf{invariantly}: it functions as a principle. When the equation
fails at some $\beta \in S(c)$ and the claimant accepts the failure, the claim is applied
\textbf{selectively}.

\begin{theorem}[Selective-Application Detection]
\label{thm:tt-selective}
Let a claimant assert $c \in \mathcal{C}_\alpha$ as a principle with scope $S(c)$, and let $c$
hold on $D_\alpha$. If $\TT_{\alpha \to \beta}(c)$ fails on some $D_\beta$ with
$\beta \in S(c)$, and the claimant accepts the failure, then the claimant's operative commitment
is a local rule of $D_\alpha$: the asserted scope exceeds the honored scope.
\end{theorem}

\begin{proof}[Argument]
A principle is a claim together with its scope. The claimant asserts the pair $(c, S(c))$ and
honors the pair $(c, \{\alpha\})$. The honored scope excludes $\beta$ by hypothesis, so the honored
scope is a proper subset of the asserted scope. The gap between the two scopes is the selective
application. \end{proof}

A failure of invariance licenses exactly one conclusion: the claim, as held by that claimant,
functions selectively. Three further conclusions remain unavailable. First, the transform reports
nothing about the truth of $c$ on any domain: a claim can hold on every axis and still be applied
selectively by a given claimant, and a claim can fail on every axis and still be applied
invariantly. Second, the transform reports nothing about the direction of repair: invariance can
be restored by extending $c$ to $D_\beta$, by retracting $c$ on $D_\alpha$, or by shrinking the
asserted scope, and the operator ranks none of these above the others. Third, the transform
measures nothing: it introduces no quantity, no calibration, and no threshold. Its output is a
scope audit of a claimant, stated in vocabulary the manuscript already supplies.

The operator therefore functions as an invariance test. It detects whether a claim is being
applied as a principle. It establishes nothing about whether the claim deserves to be one.

\section{Use I: Inter-Axis Projection}
\label{sec:tt_inter_axis}

The first use carries a claim from one axis of oppression to another and checks invariance. The
diagnostic value sits in the asymmetry the check exposes: a claimant who accepts a claim on the
gendered axis and rejects its image on the racial axis, or accepts it on the racial axis and
rejects its image on the gendered axis, applies the claim selectively, and
Theorem~\ref{thm:tt-selective} names the selectivity precisely.

\subsection{Worked Example: The Haitian Theorem on the Gendered Axis}

Chapter~\ref{ch:contradiction} states the Haitian Theorem
(Definition~\ref{def:haitian_theorem}, Equations~\ref{eq:12.5-haitian-theorem-nonkinetic}
and~\ref{eq:12.6-haitian-theorem-kinetic}): across the 1450--2026 dataset, structural liberation
of the Out-group from the extraction kernel has been achieved exclusively through kinetic action.
The theorem is established on the racial axis, with the Haitian Revolution as its anchor case.

Chapter~\ref{ch:gendered} applies the Theodore Transform to this claim. The correspondence map
runs as follows:

\begin{center}
\renewcommand{\arraystretch}{1.3}
\begin{tabular}{>{\raggedright}p{5.4cm} >{\raggedright\arraybackslash}p{8.0cm}}
\toprule
\textbf{Source vocabulary (racial axis)} & \textbf{Image vocabulary (gendered axis)} \\
\midrule
$O_{\text{racialized}}$ & $O_{\text{gendered}}$ (Equation~\ref{eq:12.7-gendered-outgroup-set}) \\
Slave patrols and the carceral apparatus (Chapter~\ref{ch:enforcement}) &
Coverture, eugenics statutes, and criminalized pregnancy (Chapter~\ref{ch:gendered}) \\
Disarmament of $O$ (Chapter~\ref{ch:kinetic}) &
Gendered kinetic disarmament (Chapter~\ref{ch:gendered}) \\
$\text{Lethal Autonomy}(O) = 0$ (Equation~\ref{eq:6.1-lethal-autonomy-gradient}) &
$\text{Lethal Autonomy}(O_{\text{gendered}}) = 0$ \\
Kinetic action & Kinetic action (fixed point of $\sigma$) \\
\bottomrule
\end{tabular}
\end{center}

The image claim states: the gendered partition, installed and maintained through force, yields
exclusively to force or to the credible, organized threat of force. Chapter~\ref{ch:gendered}
evaluates the image claim against the gendered dataset --- the witch-hunts, the colonial
dismantlement of the Agojie, the criminalization of indigenous gender systems, and the ongoing
toll of intimate-partner violence --- and records the result: the framework's conclusion for the
gendered axis is identical to its conclusion for the racial axis, because the extraction kernel is
identical. Equation~\ref{eq:tt.2-invariance} holds for this pair of domains. The Haitian Theorem
operates as a principle across the racial and gendered axes, and the gendered application of the
theorem (withdrawal, deterrence, and the necessity of force) inherits the theorem's full force.

\subsection{A Diagnostic Failure and What It Licensed}

The manuscript also records a case in which a naive projection fails, and the failure does
diagnostic work. Chapter~\ref{ch:gendered} establishes that the intersection
$O_{\text{racialized}} \cap O_{\text{gendered}}$ is transformative: the nature of the gendered
wound changes depending on which side of the racial partition the woman occupies. A correspondence
map that treats ``women'' as a monolithic Out-group produces image claims that fail on the
racialized half of the set --- the suffocation wound of the white feminine position and the
devaluation wound of the Black feminine position diverge, and a claim calibrated to one misfires
on the other.

The failure licensed a refinement of the correspondence, and the manuscript supplies the refined
object: the bivector $e_{14}$ of the Geometric Algebra extension
(Equation~\ref{eq:ga.bivector}), a two-dimensional extraction plane qualitatively distinct from
either axis alone. This is the correct reading of a failed invariance check under a coarse
correspondence. The failure localizes the point where the map requires higher resolution, and the
Geometric Algebra appendix supplies that resolution. What the failure could never license, under
Theorem~\ref{thm:tt-selective}, is continued assertion of the coarse claim at its original scope
alongside acceptance of its failure at the intersection.

\section{Use II: Cross-Domain Transposition}
\label{sec:tt_cross_domain}

The second use carries an entire formal structure between substrates. The correspondence map runs
from social mechanics to electrodynamics:

\begin{center}
\renewcommand{\arraystretch}{1.3}
\begin{tabular}{>{\raggedright}p{6.2cm} >{\raggedright\arraybackslash}p{7.2cm}}
\toprule
\textbf{Social-mechanical vocabulary} & \textbf{Electrodynamic image} \\
\midrule
Kinetic labor of $O_{\text{racialized}}$ and $I_{\text{buffer}}$ & Power supply $V_{cc}$ \\
Suppression signal injected downward from $E$ & Incident wave $a$ \\
Refusal, backlash, non-compliance & Reflected wave $b$ \\
Extraction $\mathcal{E}(t)$ & Absorbed real power (Equation~\ref{eq:xc.3-absorbed-power}) \\
Psychological wage $\psi_s$ & Reactive power $Q_{\text{reactive}}$ \\
Enclosure Score $\mathcal{S}_{\text{enc}}$ & Match condition $1 - |\Gamma|$
(Equation~\ref{eq:xc.5-enclosure-match}) \\
Cultural field $\vec{B}$ & Static magnetic bias (Theorem~\ref{thm:cultural-bias}) \\
\bottomrule
\end{tabular}
\end{center}

The derivation of Chapters~\ref{ch:redefining}--\ref{ch:algorithmic_epoch}, consolidated in
Appendix~\ref{apx:extraction_chart}, is the image of the framework under this transform. The
manuscript's electrodynamic derivation is itself an instance of the Theodore Transform, and this
instance supplies the strongest available evidence that the operator does work. The transposed
structure yielded closed theorems in the target domain --- Buffer Matching
(Theorem~\ref{thm:buffer-matching}), Reform Monotonicity
(Theorem~\ref{thm:reform-monotonicity}), Quarter-Wave Co-optation
(Theorem~\ref{thm:quarter-wave}), and Cultural Bias (Theorem~\ref{thm:cultural-bias}) --- whose
pre-images under $\TT^{-1}$ match claims the historical chapters had already established
on descriptive evidence. A mapping that manufactured agreement would produce theorems with no
source-domain correlates. The electrodynamic image produced four theorems that landed on claims
already in the record.

\subsection{Relation to the Universality Conjecture}

Appendix~\ref{app:universality} conjectures that the extraction topology is substrate-independent:
any system meeting the three conditions of Conjecture~\ref{conj:universality} converges to the
five-node topology and the phase-loading control law. The universality conjecture and the Theodore
Transform are the same argument seen from opposite ends. The conjecture asserts that shared
structure exists across substrates. The transform is the operation that exhibits the shared
structure, by carrying claims along the correspondence and checking whether they survive.
Substrate independence is the hypothesis; Equation~\ref{eq:tt.2-invariance} is the test.

\section{Falsification Criteria}
\label{sec:tt_falsification}

The claims of this appendix fail under the following conditions.

\begin{enumerate}
    \item \textbf{Selective-Application Detection.} A documented case in which the transform flags
    selective application and the flagged asymmetry is fully accounted for by a scope condition
    already stated in the framework's own vocabulary, in the manner of the scope conditions
    Chapter~\ref{ch:contradiction} states for the Haitian Theorem. A single such case marks a
    false positive; a pattern of false positives across the manuscript's own projections falsifies
    the diagnostic's discriminating power.

    \item \textbf{Correspondence failure.} A structural element the framework treats as
    load-bearing on one axis with no definable correspondent on another axis inside the asserted
    scope. If the three modes of the Tri-Modal Enclosure Model admit no statement on the gendered
    axis, for instance, the inter-axis transform is undefined for that pair of domains, and every
    invariance result depending on that pair loses standing.

    \item \textbf{Image--record contradiction.} A theorem derivable in the electrodynamic image
    whose pre-image under $\TT^{-1}$ contradicts an established result of the historical
    chapters. The transport must preserve the relations the framework asserts, and a contradiction
    between image and record is a relation destroyed. A single confirmed contradiction falsifies
    the cross-domain correspondence of Section~\ref{sec:tt_cross_domain}.

    \item \textbf{Universality linkage.} Falsification of Conjecture~\ref{conj:universality} under
    any of the three conditions listed in Appendix~\ref{app:universality} removes the structural
    ground for treating cross-domain transposition as structure-preserving, and demotes the
    electrodynamic results of this book to unconnected analogy.
\end{enumerate}

\section{Open Problems}
\label{sec:tt_open}

The transform as defined here substitutes along a hand-built correspondence. No canonical
construction of $\sigma_{\alpha\beta}$ exists, and the evidential burden of each transport rests
on the correspondence the analyst supplies. A canonical construction --- one derived from the
extraction topology itself --- would convert each application of the transform from an argued
analogy into a computed image.

Composition is unestablished. Whether
$\TT_{\beta \to \gamma} \circ \TT_{\alpha \to \beta} =
\TT_{\alpha \to \gamma}$ holds in general determines whether chains of transports
preserve structure or accumulate distortion. Until this is settled, every multi-hop projection
requires individual verification.

Completeness is uncharacterized. The set of claims invariant under all inter-axis projections has
no known decision procedure. The Haitian Theorem belongs to this set, as
Section~\ref{sec:tt_inter_axis} documents. The manuscript supplies no general method for deciding
membership for an arbitrary claim.
